% !TEX root = ../../hw01.tex
% Source: ma-gy7033fall2026homework1.pdf, p. 2, Problem 6(a--c).
% Include assignment-wide field and polynomial conventions for shared reuse.
\begin{exercise}[A polynomial determined by three values]
  \label{ex:hw01:polynomial-interpolation}
  Let \(\mathbb F\) denote \(\R\) or \(\C\), and work in
  \(\mathcal P_2(\mathbb F)\), the polynomials of degree at most two
  with coefficients in \(\mathbb F\), including the zero polynomial.
  \begin{enumerate}[(a)][2]
    \item Construct \(q_{-1},q_0,q_1\) such that, for
      \(j,k\in\{-1,0,1\}\),
      \[
        q_j(k)=
        \begin{cases}
          1,&j=k,\\
          0,&j\ne k.
        \end{cases}
      \]
    \item Prove that these polynomials form a basis and that every
      \(p\in\mathcal P_2(\mathbb F)\) satisfies
      \[
        p=p(-1)q_{-1}+p(0)q_0+p(1)q_1.
      \]
    \item Find the unique polynomial \(p\in\mathcal P_2(\mathbb F)\)
      for which \(p(-1)=2\), \(p(0)=-1\), and \(p(1)=4\).
  \end{enumerate}
  \solutionlink{hw01:polynomial-interpolation}
\end{exercise}
